% ==================================================================
% id:       lim-036
% topic:    גבולות ורציפות של פונקציות
% subtopic: אריתמטיקה של גבולות
% solution: full
% ==================================================================
\textbf{הוכיחו או הפריכו כל אחד מהסעיפים הבאים:}
\begin{enumerate}
\item אם קיימים הגבולות $\ds\lim_{x\to x_0}f(x)$ ו-$\ds\lim_{x\to x_0}\lr{f(x)+g(x)}$, אז קיים הגבול $\ds\lim_{x\to x_0}g(x)$.
\ifsol
\hint{כתבו את $g$ בעזרת $f+g$ ו-$f$.}
\begin{solbox}
הטענה נכונה. מתקיים $g(x)=\lr{f(x)+g(x)}-f(x)$, ולפי אריתמטיקה של גבולות (הפרש של שתי פונקציות שיש להן גבול) קיים הגבול של $g$ ב-$x_0$.
\end{solbox}
\fi

    \item
    אם
     קיים הגבול
    $\displaystyle\lim_{x\to x_0}\left(f(x)+g(x)\right)$,
    אזי קיימים הגבולות
    $\displaystyle\lim_{x\to x_0} f(x)$
    ו-
    $\displaystyle\lim_{x\to x_0} g(x)$.
\ifsol
\hint{אם נראה לכם שהטענה נכונה, נסו להביע את $f$ בעזרת $f+g$ ו-$g$. אחרת, חפשו דוגמה נגדית פשוטה.}
\begin{solbox}
הטענה אינה נכונה.
תהיינה
$$f(x)=\frac{1}{x},\quad
g(x)=-f(x)$$
ו-
$x_0=0$.
מתקיים
$\displaystyle\lim_{x\to0}(f(x)+g(x))=0$
כי
$f(x)+g(x)=0$
לכל
$x\neq0$, ומצד שני
בבירור לא קיימים הגבולות
$\displaystyle\lim_{x\to0}f(x)$
ו-
$\displaystyle\lim_{x\to0}(g(x))$.

שימו לב כי כל עוד
$\ds{\lim_{x\to x_0}}f(x)$
לא קיים, הדוגמא
$g=-f$
"תעבוד".
\end{solbox}
\fi
\item אם קיימים הגבולות $\ds\lim_{x\to x_0}f(x)$ ו-$\ds\lim_{x\to x_0}f(x)g(x)$, אז קיים הגבול $\ds\lim_{x\to x_0}g(x)$.
\ifsol
\begin{solbox}
הטענה אינה נכונה. למשל $x_0=0$, $f(x)=0$ ו-$g(x)=\frac1x$. אז $f(x)g(x)=0$ לכל $x\ne0$, ולכן לשתי הפונקציות $f$ ו-$fg$ יש גבול $0$. אבל ל-$\frac1x$ אין גבול ב-$0$: עבור $a_n=\frac1n$ מתקיים $g(a_n)=n\to\infty$.
\end{solbox}
\fi
\item אם $\ds\lim_{x\to x_0}f(x)=2$ ו-$\ds\lim_{x\to x_0}f(x)g(x)=6$, אז $\ds\lim_{x\to x_0}g(x)=3$.
\ifsol
\hint{האם מותר לחלק ב-$f$ ליד $x_0$?}
\begin{solbox}
הטענה נכונה. מכיוון שהגבול של $f$ שונה מ-$0$, בסביבה נקובה של $x_0$ מתקיים $f(x)\ne0$, ושם
\[g(x)=\frac{f(x)g(x)}{f(x)}\]
ולכן
\[\lim_{x\to x_0}g(x)=\lim_{x\to x_0}\frac{f(x)g(x)}{f(x)}=\frac62=3\]
לפי אריתמטיקה של גבולות.
\end{solbox}
\fi
\item אם ל-$f$ אין גבול ב-$x_0$, אז ל-$f^2$ אין גבול ב-$x_0$.
\ifsol
\begin{solbox}
הטענה אינה נכונה. למשל $x_0=0$ ו-
\[f(x)=\begin{cases}1, & x\ge0\\ -1, & x<0.\end{cases}\]
לפי היינה ל-$f$ אין גבול ב-$0$: עבור $\frac1n$ מתקבל $1$, ועבור $-\frac1n$ מתקבל $(-1)$. אבל $f^2(x)=1$ לכל $x$, ולכן $\ds\lim_{x\to0}f^2(x)=1$.
\end{solbox}
\fi
\item אם $\ds\lim_{x\to x_0}g(x)=0$ וקיים הגבול (הסופי) $\ds\lim_{x\to x_0}\frac{f(x)}{g(x)}$, אז $\ds\lim_{x\to x_0}f(x)=0$.
\ifsol
\begin{solbox}
הטענה נכונה. נסמן $\ds\lim_{x\to x_0}\frac{f(x)}{g(x)}=L$. בסביבה נקובה של $x_0$ המנה מוגדרת, כלומר $g(x)\ne0$, ושם
\[f(x)=\frac{f(x)}{g(x)}\cdot g(x)\]
ולכן
\[\lim_{x\to x_0}f(x)=\lim_{x\to x_0}\left(\frac{f(x)}{g(x)}\cdot g(x)\right)=L\cdot0=0\]
\end{solbox}
\fi
\end{enumerate}
